\chapter[Równanie falowe]{Równanie falowe: charakterystyki, energia i propagacja}
\label{chap:pde-wave}

Równanie ciepła wygładza ograniczone dane początkowe,
a niezerowe dane nieujemne o zwartym nośniku prowadzą
do rozwiązania dodatniego w każdym punkcie przestrzeni
w każdej dodatniej chwili. Równanie falowe
zachowuje się inaczej: zaburzenie nie przekracza określonej
prędkości, a oscylacje nie są samoczynnie tłumione. W tym
rozdziale wyprowadzimy jawne wzory, udowodnimy jednoznaczność
na podstawie energii i rozróżnimy słabą oraz silną zasadę
Huygensa.

Ustalmy całkowite $n\geq1$, $c>0$ i oznaczmy
$\square_c u=\partial_t^2u-c^2\Delta_xu$. Zagadnienie
Cauchy'ego ma postać
\begin{equation}\label{eq:pde5-wave}
 \square_cu=0\quad\text{dla }t>0,\ x\in\R^n,\qquad
 u(0,x)=g(x),\quad u_t(0,x)=h(x).
\end{equation}
Dwa warunki początkowe są potrzebne, ponieważ równanie
jest drugiego rzędu względem czasu. Funkcje skalarne
są w tym rozdziale rzeczywiste, o ile nie zaznaczono
inaczej. Przez $\lambda_k$ oznaczamy miarę Lebesgue'a
w $\R^k$ dla $k\geq1$, a przez $\sigma$ standardową
miarę powierzchniową na rozważanej gładkiej
hiperpowierzchni; w wymiarze jeden na brzegu
przedziału jest to miara licząca.
Przez $C^2$ na domkniętym pasku czasowym rozumiemy
funkcje o ciągłych pochodnych do rzędu drugiego
także na jego końcach, gdzie pochodne czasowe
są jednostronne. Wszystkie
rachunki z całkowaniem i różniczkowaniem wykonujemy
najpierw dla danych dostatecznie gładkich; słabe
rozwiązania wprowadzimy dopiero w kolejnych częściach.

\section{Równanie w jednym wymiarze: rozkład na fale biegnące}
\label{sec:pde5-dalembert}

Rozważmy $u_{tt}-c^2u_{xx}=0$ dla $(t,x)\in(0,\infty)\times\R$.
Operator rozkłada się na czynniki
$\partial_t^2-c^2\partial_x^2
=(\partial_t-c\partial_x)(\partial_t+c\partial_x)$.
Wprowadzamy współrzędne charakterystyczne
$\xi=x-ct$, $\eta=x+ct$. Odwzorowanie
$(t,x)\mapsto(\xi,\eta)$ jest odwracalne, a jego
odwrotność dana jest wzorami
$t=(\eta-\xi)/(2c)$ oraz $x=(\eta+\xi)/2$.
Półpłaszczyźnie $t>0$ odpowiada zbiór
$D=\{(\xi,\eta)\in\R^2:\xi<\eta\}$.
Jeżeli $U(\xi,\eta)=u((\eta-\xi)/(2c),(\eta+\xi)/2)$,
to z reguły łańcuchowej
\[
 \partial_tu=-cU_\xi+cU_\eta,\qquad
 \partial_xu=U_\xi+U_\eta,\qquad
 u_{tt}-c^2u_{xx}=-4c^2U_{\xi\eta}.
\]
W ostatniej równości używamy przemienności
pochodnych mieszanych dla funkcji klasy $C^2$.

Wykażmy globalną, a nie tylko lokalną reprezentację
rozwiązania na $D$. Z $U_{\xi\eta}=0$ wynika, że
dla każdego ustalonego $\xi$ funkcja
$\eta\mapsto U_\xi(\xi,\eta)$ jest stała na spójnym
przedziale $(\xi,\infty)$. Podobnie dla ustalonego
$\eta$ funkcja $\xi\mapsto U_\eta(\xi,\eta)$
jest stała na $(-\infty,\eta)$.
Definiujemy $a(\xi)=U_\xi(\xi,\xi+1)$ oraz
$b(\eta)=U_\eta(\eta-1,\eta)$. Funkcje $a,b$
są klasy $C^1$ na $\R$ i na $D$ zachodzi
$U_\xi(\xi,\eta)=a(\xi)$, $U_\eta(\xi,\eta)=b(\eta)$.
Wybierając funkcje pierwotne $F'=a$ i $G'=b$,
otrzymujemy $\nabla(U-F-G)=0$ na $D$, gdzie
$F=F(\xi)$ i $G=G(\eta)$. Zbiór $D$ jest wypukły,
więc różnica jest stała: jej pochodna wzdłuż
dowolnego odcinka w $D$ znika.
Włączając tę stałą do $F$, dostajemy
\[
 u(t,x)=F(x-ct)+G(x+ct),\qquad F,G\in C^2(\R).
\]
Jeżeli $u$ jest klasy $C^2$ także przy $t=0$,
reprezentacja i wzory na pochodne przedłużają się
do $t=0$ przez ciągłość. Ten argument uzasadnia
reprezentację na całej półpłaszczyźnie, potrzebną
w dowodzie jednoznaczności.

W rozdziale wprowadzającym otrzymaliśmy już wzór d'Alemberta
(twierdzenie~\ref{thm:dalembert}). Wracamy do niego tutaj,
aby powiązać reprezentację z charakterystykami, energią i
obszarem zależności. Powtórzony krótki dowód ma charakter
samodzielnego przypomnienia, nie stanowi odrębnego twierdzenia
o szerszych założeniach.

\begin{twierdzenie}[Wzór d'Alemberta]\label{thm:pde5-dalembert}
Niech $c>0$, $g\in C^2(\R)$ i $h\in C^1(\R)$.
Istnieje dokładnie jedno rozwiązanie
$u\in C^2([0,\infty)\times\R)$ zagadnienia
\eqref{eq:pde5-wave} dla $n=1$. Jest nim
\begin{equation}\label{eq:pde5-dalembert}
 u(t,x)=\frac{g(x-ct)+g(x+ct)}2
 +\frac1{2c}\int_{x-ct}^{x+ct}h(y)\,\mathrm d\lambda_1(y).
\end{equation}
\end{twierdzenie}
\begin{proof}
Najpierw wyprowadzimy konieczną postać rozwiązania.
Na mocy globalnej reprezentacji powyżej każde
rozwiązanie ma postać $F(x-ct)+G(x+ct)$.
Warunki początkowe i ciągłość pochodnych przy $t=0$
dają $F+G=g$ oraz $-cF'+cG'=h$ na $\R$.
Różniczkując pierwszą równość i rozwiązując układ,
dostajemy $F'=(g'-h/c)/2$ oraz $G'=(g'+h/c)/2$.
Niech $H(x)=\int_0^x h(y)\,\mathrm d\lambda_1(y)$,
gdzie dla $x<0$ całkę rozumiemy z przeciwną orientacją.
Ponieważ $h\in C^1(\R)$, mamy $H\in C^2(\R)$
oraz $H'=h$. Zatem
\[
 F(x)=\tfrac12g(x)-\tfrac1{2c}H(x)+A,\qquad
 G(x)=\tfrac12g(x)+\tfrac1{2c}H(x)+B.
\]
Warunek $F+G=g$ wymusza $A+B=0$.
Po wstawieniu argumentów $x-ct$ i $x+ct$
otrzymujemy~\eqref{eq:pde5-dalembert}.
W szczególności każde rozwiązanie musi mieć
tę samą postać, co daje jednoznaczność.

Pozostaje wykazać istnienie. Wzór po prawej stronie
\eqref{eq:pde5-dalembert} definiuje funkcję klasy
$C^2$ (nawet dla wszystkich $t\in\R$, przy
orientowanej całce). Oznaczmy $a=x-ct$, $b=x+ct$.
Z reguły łańcuchowej i twierdzenia podstawowego
rachunku całkowego otrzymujemy
\[
 \begin{aligned}
 u_t&=\tfrac c2(-g'(a)+g'(b))+\tfrac12(h(a)+h(b)),\\
 u_x&=\tfrac12(g'(a)+g'(b))+\tfrac1{2c}(h(b)-h(a)),\\
 u_{tt}&=\tfrac{c^2}2(g''(a)+g''(b))
                   +\tfrac c2(h'(b)-h'(a)),\\
 u_{xx}&=\tfrac12(g''(a)+g''(b))
                   +\tfrac1{2c}(h'(b)-h'(a)).
 \end{aligned}
\]
Stąd $u_{tt}=c^2u_{xx}$. W chwili $t=0$ mamy
$a=b=x$, więc $u(0,x)=g(x)$ oraz $u_t(0,x)=h(x)$.
Wzór istotnie daje rozwiązanie zagadnienia Cauchy'ego.
Dowód nie wymaga ograniczoności ani kontroli wzrostu
danych lub rozwiązania w nieskończoności.
\end{proof}

\begin{uwaga}[Dokładny obszar zależności]
Ze wzoru wynika, że wartość $u(t,x)$ zależy jedynie
od $g$ w dwóch punktach $x-ct,x+ct$ i od $h$ na
odcinku $[x-ct,x+ct]$. W szczególności dane z punktów
poza tym odcinkiem nie wpływają na $u(t,x)$.
Dokładniej, dla $t>0$ zmiana wartości $g$ poza
końcami odcinka nie zmienia $u(t,x)$, natomiast
zmiana $h$ o funkcję gładką, nieujemną, niezerową
i o nośniku zawartym w dowolnym niepustym
podprzedziale $(x-ct,x+ct)$ zmienia tę wartość
o dodatnią liczbę. Nie oznacza to więc, że
sygnał jest odczuwany wyłącznie na czole fali:
człon z $h$ zależy także od danych we wnętrzu
podstawy stożka wstecznego.
\end{uwaga}

\begin{przyklad}[Dwa przesuwające się impulsy]
Jeżeli $h=0$ i $g\in C_c^2(\R)$, to
$u(t,x)=\frac12g(x-ct)+\frac12g(x+ct)$.
Początkowy impuls rozdziela się na dwa profile
biegnące w przeciwnych kierunkach bez zmiany kształtu.
Jeżeli natomiast $g=0$, $h\in C_c^1(\R)$,
$h\geq0$ i $h\not\equiv0$, to $\int_\R h\,\mathrm d\lambda_1>0$:
z ciągłości funkcja jest dodatnia i ograniczona
z dołu przez dodatnią stałą na pewnym przedziale.
Wybierzmy $R>0$ takie, że $\operatorname{supp}h\subset[-R,R]$.
Dla ustalonego $x$ oraz $t\geq(|x|+R)/c$ mamy
$[-R,R]\subset[x-ct,x+ct]$, a więc
$u(t,x)=(2c)^{-1}\int_\R h(y)\,\mathrm d\lambda_1(y)>0$.
Ta druga sytuacja pokazuje, dlaczego silna zasada
Huygensa nie zachodzi w wymiarze jeden.
\end{przyklad}

\begin{center}
\begin{tikzpicture}[scale=.88,>=Stealth]
  \fill[blue!8] (-2.7,0)--(2.7,0)--(0,3.35)--cycle;
  \draw[->] (-4,0)--(4,0) node[right] {$x$};
  \draw[->] (0,-.25)--(0,3.9) node[above] {$t$};
  \draw[blue!65!black,thick] (-2.7,0)--(0,3.35)--(2.7,0);
  \fill (0,3.35) circle (1.8pt) node[above right] {$(t_0,x_0)$};
  \draw[<->,blue!50!black] (-2.7,-.3)--(2.7,-.3)
    node[midway,below] {$[x_0-ct_0,x_0+ct_0]$};
  \node[blue!65!black] at (0,1.05) {stożek wsteczny};
  \node at (3.4,3.2) {$x_0=0$};
\end{tikzpicture}
\par\smallskip\small
Obszar zależności w jednym wymiarze; osie i wierzchołek
narysowano dla $x_0=0$, a szerokość stożka odpowiada $c t_0$.
\end{center}

\section{Energia i jednoznaczność w wielu wymiarach}
\label{sec:pde5-energy}

\begin{definicja}[Energia falowa]
Niech $T>0$, $c>0$ i $u\in C^2([0,T]\times\R^n)$.
Jeżeli dla ustalonego $t\in[0,T]$ pochodne
$u_t(t,\cdot)$ oraz $\nabla_xu(t,\cdot)$ są
całkowalne z kwadratem, definiujemy energię
w tej chwili przez
\begin{equation}\label{eq:pde5-energy}
 E_u(t)=\frac12\int_{\R^n}
 \bigl(|u_t(t,x)|^2+c^2|\nabla_xu(t,x)|^2\bigr)
 \,\mathrm d\lambda_n(x).
\end{equation}
\end{definicja}

\begin{twierdzenie}[Zachowanie energii]\label{thm:pde5-energy}
Niech $c,T>0$, $n\ge1$, a $u\in C^2([0,T]\times\R^n)$
spełnia $\square_cu=0$. Załóżmy, że istnieje
zwarty zbiór $K\subset\R^n$, poza którym
$u_t(t,\cdot)=\nabla_xu(t,\cdot)=0$
dla każdego $t\in[0,T]$. Wówczas $E_u(t)=E_u(0)$
dla $0\leq t\leq T$.
\end{twierdzenie}
\begin{proof}
Wybierzmy $R>0$ tak, aby $K\subset B(0,R)$,
i oznaczmy $B=B(0,R)$.
Dla każdego $t\in[0,T]$ całkę definiującą energię
można ograniczyć do $B$, ponieważ jej funkcja
podcałkowa znika poza $K$.
Wszystkie pochodne występujące w różniczkowaniu
tej funkcji względem czasu są ciągłe i ograniczone
na zwartym zbiorze $[0,T]\times\overline B$.
Z twierdzenia o różniczkowaniu pod całką,
uzasadnionego przez zbieżność zdominowaną na
zbiorze skończonej miary, wynika
\[
 E_u'(t)=\int_B
 \bigl(u_tu_{tt}+c^2\nabla_xu\cdot\nabla_xu_t\bigr)
 \,\mathrm d\lambda_n(x).
\]
W szczególności energia jest klasy $C^1$ względem
czasu, z pochodnymi jednostronnymi na końcach.
Z równości
$\diver(u_t\nabla_xu)=\nabla_xu_t\cdot\nabla_xu+
u_t\Delta_xu$ oraz twierdzenia Gaussa otrzymujemy
\[
 \begin{aligned}
 E_u'(t)
 &=\int_B u_t(u_{tt}-c^2\Delta_xu)
      \,\mathrm d\lambda_n(x)\\
 &\quad+c^2\int_{\partial B}u_t\partial_\nu u
      \,\mathrm d\sigma=0.
 \end{aligned}
\]
Pierwsza całka znika na mocy równania, a druga
dlatego, że obie pochodne $u_t$ i $\nabla_xu$
znikają na $\partial B$.
Dla $n=1$ twierdzenie Gaussa zastępujemy
twierdzeniem podstawowym rachunku całkowego na
przedziale $(-R,R)$; wartości brzegowe strumienia
również znikają. Stąd $E_u(t)=E_u(0)$ dla każdego
$t\in[0,T]$.
Założenie o wspólnym zwartym nośniku jest tutaj
warunkiem wystarczającym do przeprowadzenia dowodu,
nie warunkiem koniecznym zachowania energii.
Sama skończoność energii w poszczególnych chwilach
nie uzasadnia powyższego różniczkowania pod całką.
\end{proof}

Najważniejszy wynik nie wymaga jednak założeń o zachowaniu
rozwiązania przy $|x|\to\infty$. Zamiast całej przestrzeni
rozważymy kurczące się kule.

\begin{twierdzenie}[Lokalna energia i skończona prędkość propagacji]
\label{thm:pde5-cone}
Niech $n\ge1$, $c,T>0$, $x_0\in\R^n$ oraz
$u\in C^2([0,T]\times\R^n)$ spełnia $\square_cu=0$.
Jeżeli $u(0,x)=u_t(0,x)=0$ dla wszystkich
$x\in\overline B(x_0,cT)$, to $u(T,x_0)=0$.
Co więcej, w stożku
$\{(t,x):0\le t<T,\ |x-x_0|<c(T-t)\}$
mamy $u(t,x)=0$.
\end{twierdzenie}
\begin{proof}
Dla $0\le t<T$ połóżmy $r(t)=c(T-t)$,
$B_t=B(x_0,r(t))$ i
$e=\frac12(u_t^2+c^2|\nabla_xu|^2)$. Definiujemy
$E(t)=\int_{B_t}e(t,x)\,\mathrm d\lambda_n(x)$.
Dla $n\geq2$ zapisujemy całkę we współrzędnych
biegunowych o środku $x_0$:
\[
 E(t)=\int_0^{r(t)}\!\int_{S^{n-1}}
 e(t,x_0+\rho\theta)\rho^{n-1}\,
 \mathrm d\sigma(\theta)\,\mathrm d\lambda_1(\rho).
\]
Funkcja $e$ i jej pochodna czasowa są ciągłe
na zwartym zbiorze
$[0,T]\times\overline B(x_0,cT)$.
Wewnętrzną całkę możemy więc różniczkować
względem $t$ pod znakiem całki, a otrzymana
pochodna jest ciągła względem $(t,\rho)$
dla $0\leq\rho\leq cT$. Reguła Leibniza dla całki
o zmiennej górnej granicy daje
\[
 \begin{aligned}
 E'(t)
 &=\int_{B_t}e_t(t,x)\,\mathrm d\lambda_n(x)\\
 &\quad+r'(t)r(t)^{n-1}
     \int_{S^{n-1}}e(t,x_0+r(t)\theta)
       \,\mathrm d\sigma(\theta)\\
 &=\int_{B_t}e_t(t,x)\,\mathrm d\lambda_n(x)
    -c\int_{\partial B_t}e(t,x)\,\mathrm d\sigma(x).
 \end{aligned}
\]
W ostatnim przejściu użyliśmy skalowania miary
powierzchniowej przez czynnik $r(t)^{n-1}$
oraz $r'(t)=-c$.
Dla $n=1$ stosujemy bezpośrednio regułę Leibniza
do całki od $x_0-r(t)$ do $x_0+r(t)$: oba
końce dają składniki $-ce(t,x_0\pm r(t))$.
W obu przypadkach $E$ jest klasy $C^1$ dla
$0\leq t<T$, przy pochodnej prawostronnej w zerze.

Z równania falowego wynika punktowa tożsamość
\[
 e_t=u_tu_{tt}+c^2\nabla_xu\cdot\nabla_xu_t
     =c^2\diver(u_t\nabla_xu).
\]
Pole $u_t\nabla_xu$ jest klasy $C^1$ względem $x$,
więc twierdzenie Gaussa na kuli (w wymiarze jeden
twierdzenie podstawowe rachunku całkowego) daje
\[
\begin{aligned}
 E'(t)&=c^2\int_{\partial B_t}u_t\partial_\nu u
 \,\mathrm d\sigma-c\int_{\partial B_t}e\,\mathrm d\sigma\\
 &=-\frac c2\int_{\partial B_t}
 \bigl((u_t-c\partial_\nu u)^2+c^2|\nabla_{\tan}u|^2\bigr)
 \,\mathrm d\sigma\leq0.
\end{aligned}
\]
Tutaj $\nu$ jest zewnętrzną normalną jednostkową,
$\partial_\nu u=\nabla_xu\cdot\nu$, zaś
$\nabla_{\tan}u=\nabla_xu-(\partial_\nu u)\nu$.
Dla $n=1$ rozumiemy całkę brzegową jako sumę po
dwóch końcach przedziału, a składnik styczny jest
zerowy. Do dopełnienia kwadratu użyliśmy
ortogonalnego rozkładu gradientu na składową
normalną i styczną:
$|\nabla_xu|^2=|\partial_\nu u|^2+|\nabla_{\tan}u|^2$.
Z założeń początkowych $u(0,\cdot)$ jest
identycznie zerowa na otwartej kuli $B_0$,
więc także jej gradient znika na tej kuli.
Wraz z $u_t(0,\cdot)=0$ daje to $E(0)=0$.
Ponieważ $E'\leq0$ oraz $E\geq0$,
mamy $E(t)=0$ dla $0\leq t<T$.
Dla ustalonego $t$ nieujemna, ciągła funkcja
$e(t,\cdot)$ musi zatem znikać w każdym punkcie
$B_t$: gdyby była dodatnia w pewnym punkcie,
byłaby ograniczona z dołu przez dodatnią stałą
na pewnej kuli dodatniej miary, wbrew $E(t)=0$.
W szczególności $u_t=\nabla_xu=0$ w rozważanym
stożku. Jeżeli $x\in B_t$, to $x\in B_s$ dla
wszystkich $0\leq s\leq t$, gdyż promień kul
maleje. Stąd
$u(t,x)=u(0,x)+\int_0^t u_t(s,x)\,
\mathrm d\lambda_1(s)=0$.
Wreszcie $x_0\in B_t$ dla wszystkich $t<T$,
więc z ciągłości $u(T,x_0)=0$.
\end{proof}


\begin{twierdzenie}[Lokalne oszacowanie energii dla wymuszonego równania falowego]
\label{thm:pde5-cone-forcing}
Niech $n\ge1$, $c,T>0$, $x_0\in\R^n$, a
$u\in C^2([0,T]\times\R^n)$ spełnia
$u_{tt}-c^2\Delta_xu=f$ dla ciągłej funkcji
$f:[0,T]\times\R^n\to\R$.
Dla $0\le t<T$ oznaczmy $B_t=B(x_0,c(T-t))$ i
\[
 E(t)=\frac12\int_{B_t}
 \bigl(|u_t(t,x)|^2+c^2|\nabla_xu(t,x)|^2\bigr)
 \,\mathrm d\lambda_n(x).
\]
Wtedy dla każdego $0\le t<T$ zachodzi
\begin{equation}\label{eq:pde5-forced-cone-bound}
 \sqrt{E(t)}\le\sqrt{E(0)}+
 \frac1{\sqrt2}\int_0^t
 \|f(s,\cdot)\|_{L^2(B_s)}\,\mathrm d\lambda_1(s).
\end{equation}
W szczególności jeśli początkowe dane $u,u_t$ znikają
na $\overline B(x_0,cT)$, zaś $f=0$ w stożku
$\{(t,x):0\le t<T,\ x\in B_t\}$, to $u$ znika
w tym stożku oraz w jego wierzchołku $(T,x_0)$.
\end{twierdzenie}
\begin{proof}
Z twierdzenia o pochodnej całki po ruchomej kuli,
przeprowadzonego w dowodzie twierdzenia~\ref{thm:pde5-cone},
oraz z tożsamości
$\partial_t e=c^2\diver(u_t\nabla_xu)+u_tf$
otrzymujemy
\[
 E'(t)=-\frac c2\int_{\partial B_t}
 \bigl((u_t-c\partial_\nu u)^2+
 c^2|\nabla_{\tan}u|^2\bigr)\,\mathrm d\sigma
 +\int_{B_t}u_t(t,x)f(t,x)\,\mathrm d\lambda_n(x).
\]
Pierwszy składnik jest niedodatni. Dla drugiego
z nierówności Cauchy'ego--Schwarza otrzymujemy
\[
 E'(t)\leq
 \|u_t(t,\cdot)\|_{L^2(B_t)}\|f(t,\cdot)\|_{L^2(B_t)}
 \leq\sqrt{2E(t)}\|f(t,\cdot)\|_{L^2(B_t)}.
\]
Norma wymuszenia na ruchomej kuli jest ciągła
względem czasu dla $t<T$: po podstawieniu
$x=x_0+r(t)y$ jej kwadrat jest równy
$r(t)^n\int_{B(0,1)}|f(t,x_0+r(t)y)|^2
\,\mathrm d\lambda_n(y)$, a ciągłość wynika
ze zbieżności zdominowanej na ustalonej kuli,
gdyż $f$ jest ograniczona na domkniętym stożku.
Nie wolno bez uzasadnienia dzielić przez $\sqrt{E(t)}$,
ponieważ energia może znikać. Dla $\delta>0$
z reguły łańcuchowej i powyższej nierówności mamy
\[
 \begin{aligned}
 \frac{\mathrm d}{\mathrm dt}\sqrt{E(t)+\delta}
 &=\frac{E'(t)}{2\sqrt{E(t)+\delta}}\\
 &\leq\frac1{\sqrt2}\sqrt{\frac{E(t)}{E(t)+\delta}}\,
       \|f(t,\cdot)\|_{L^2(B_t)}\\
 &\leq\frac1{\sqrt2}\|f(t,\cdot)\|_{L^2(B_t)}.
 \end{aligned}
\]
Całkując na $[0,t]$, otrzymujemy oszacowanie
$\sqrt{E(t)+\delta}\leq\sqrt{E(0)+\delta}+
2^{-1/2}\int_0^t\|f(s,\cdot)\|_{L^2(B_s)}
\,\mathrm d\lambda_1(s)$.
Prawa strona poza wyrazem początkowym nie zależy
od $\delta$. Przechodząc z $\delta\downarrow0$,
dostajemy~\eqref{eq:pde5-forced-cone-bound}.
Jeżeli dane początkowe znikają na $\overline B_0$,
to również gradient danych $u(0,\cdot)$ znika na
$B_0$, zatem $E(0)=0$. Jeśli ponadto $f=0$
w rozważanym stożku, prawa strona oszacowania
jest zerowa dla każdego $t<T$. Argument z
dowodu twierdzenia~\ref{thm:pde5-cone} daje
punktowe zanikanie pochodnych, a następnie
$u(t,x)=0$ w stożku przez całkowanie po pionowym
odcinku zawartym w stożku. Ciągłość daje
$u(T,x_0)=0$.
\end{proof}

\begin{wniosek}[Jednoznaczność bez ograniczeń wzrostu]
\label{cor:pde5-uniqueness}
Niech $n\ge1$, $c,T>0$ i niech
$u,v\in C^2([0,T]\times\R^n)$ rozwiązują to samo
jednorodne równanie falowe i mają takie same dane
początkowe w całej $\R^n$. Wówczas $u=v$.
\end{wniosek}
\begin{proof}
Do $w=u-v$ stosujemy twierdzenie~\ref{thm:pde5-cone}
dla dowolnego $(T',x_0)$, $0<T'\le T$;
ograniczenie $w$ do $[0,T']\times\R^n$ ma
wymaganą regularność, spełnia równanie i ma
zerowe dane początkowe na $\overline B(x_0,cT')$.
Stąd $w(T',x_0)=0$. Ponieważ $T'$ i $x_0$
były dowolne, a w chwili $0$ dane są równe,
otrzymujemy $w=0$ na całym pasku.

W ten sam sposób stosując
twierdzenie~\ref{thm:pde5-cone-forcing} do różnicy
dwóch rozwiązań niejednorodnych, otrzymujemy
lokalną zasadę zależności od danych:
jeżeli oba profile początkowe są zgodne na
$\overline B(x_0,cT')$ i wymuszenia są zgodne
w stożku wstecznym punktu $(T',x_0)$, to wartości
rozwiązań w tym punkcie są równe.
Nie potrzeba zakładać skończonej energii
globalnej ani ograniczoności rozwiązań.
\end{proof}

\begin{uwaga}[Kontrast z równaniem ciepła]
W rozdziale~\ref{chap:pde-heat} skonstruowaliśmy
niezerowe, gładkie rozwiązanie Tychonoffa o zerowych
danych początkowych, co wymagało porzucenia ograniczeń
wzrostu. Dla klasycznego równania falowego na całej
$\R^n$ taki fenomen nie występuje w klasie $C^2$:
lokalna tożsamość energetyczna zapewnia jednoznaczność
bez jakichkolwiek warunków przy $|x|\to\infty$.
\end{uwaga}

\section{Wzór Kirchhoffa w trzech wymiarach}
\label{sec:pde5-kirchhoff}

Dla $v\in C^2(\R^3)$, $x\in\R^3$ i $r>0$
wprowadzamy średnią sferyczną
\begin{equation}\label{eq:pde5-sphere}
 M_rv(x)=\frac1{4\pi}\int_{S^2}v(x+r\omega)
 \,\mathrm d\sigma_{S^2}(\omega),
\end{equation}
gdzie $S^2=\{\omega\in\R^3:|\omega|=1\}$,
a $\mathrm d\sigma_{S^2}$ jest miarą powierzchniową
o całkowitej masie $4\pi$. Definiujemy $M_0v(x)=v(x)$.

\begin{lemat}[Równanie dla średnich sferycznych]
\label{lem:pde5-darboux}
Jeżeli $v\in C^2(\R^3)$, to dla $r>0$
\begin{equation}\label{eq:pde5-darboux}
 \partial_r^2 M_rv(x)+\frac2r\partial_rM_rv(x)
 =M_r(\Delta v)(x).
\end{equation}
\end{lemat}
\begin{proof}
Dla ustalonego $x$ niech $V(r,\omega)=v(x+r\omega)$.
Rozkład laplasjanu na składową radialną i sferyczną daje
$\Delta v(x+r\omega)=\partial_r^2V+(2/r)\partial_rV
+r^{-2}\Delta_{S^2}V$. Całka po sferze z
$\Delta_{S^2}V$ jest zerowa, ponieważ sfera nie
ma brzegu (można to sprawdzić przez lokalną
formułę dywergencji na sferze). Całkujemy otrzymaną
tożsamość po $S^2$ i dzielimy przez $4\pi$.
\end{proof}

\begin{twierdzenie}[Wzór Kirchhoffa]
\label{thm:pde5-kirchhoff}
Niech $c>0$, $g\in C^3(\R^3)$ oraz
$h\in C^2(\R^3)$. Wówczas istnieje dokładnie
jedno rozwiązanie $u\in C^2([0,\infty)\times\R^3)$
zagadnienia~\eqref{eq:pde5-wave}, dane wzorem
\begin{equation}\label{eq:pde5-kirchhoff}
 u(t,x)=\frac{\mathrm d}{\mathrm dt}
 \bigl(t M_{ct}g(x)\bigr)+t M_{ct}h(x),\qquad t\ge0.
\end{equation}
\end{twierdzenie}
\begin{proof}
Najpierw załóżmy, że $u$ jest rozwiązaniem $C^2$.
Dla ustalonego $x$ kładziemy
$W(t,r)=rM_ru(t,\cdot)(x)$, $r\ge0$. Z lematu
\ref{lem:pde5-darboux} otrzymujemy dla $r>0$
$W_{tt}-c^2W_{rr}=0$, gdyż
$\partial_r^2(rM_ru)=rM_r\Delta u$.
Mamy $W(t,0)=0$,
$W(0,r)=rM_rg(x)$ i $W_t(0,r)=rM_rh(x)$.
Przedłużając dane nieparzyście na $r<0$
i stosując wzór d'Alemberta z twierdzenia
\ref{thm:pde5-dalembert}, otrzymujemy
$W_r(t,0)=\frac{\mathrm d}{\mathrm dr}(rM_rg)|_{r=ct}
 +c^{-1}(rM_rh)|_{r=ct}$.
Ponieważ $W_r(t,0)=u(t,x)$, dostajemy
\eqref{eq:pde5-kirchhoff}.

Dla istnienia zdefiniujmy dla $v\in C^2(\R^3)$
$V_v(t,x)=tM_{ct}v(x)$. Z tożsamości Darboux
otrzymujemy dla $t>0$ bezpośrednio
\[
 (\partial_t^2-c^2\Delta_x)V_v
 =t c^2\left(\partial_r^2M_rv+\frac2r\partial_rM_rv
              -M_r\Delta v\right)\bigg|_{r=ct}=0.
\]
Dla $g\in C^3(\R^3)$ otrzymujemy $\partial_tV_g\in C^2$,
a dla $h\in C^2(\R^3)$ także $V_h\in C^2$.
Aby uzasadnić, że $\square_c(\partial_tV_g)=0$
przy tej skończonej regularności, aproksymujemy $g$
lokalnie w $C^3$ przez funkcje gładkie; skończony
zasięg całek sferycznych daje lokalną zbieżność
$\partial_tV_{g_j}\to\partial_tV_g$ w $C^2$.
Dla danych gładkich można skomutować $\square_c$
z $\partial_t$, a przejście do granicy dowodzi żądanej równości.
Ponadto $M_0v=v$ i $\partial_rM_rv|_{r=0}=0$,
ponieważ średnia z $\omega$ na $S^2$ jest zerowa.
Zatem $u=\partial_tV_g+V_h$ spełnia
$u(0,x)=g(x)$ i $u_t(0,x)=h(x)$.
Jednoznaczność wynika z wniosku~\ref{cor:pde5-uniqueness}.
\end{proof}

\begin{uwaga}[Jakie dane widzi punkt w trzech wymiarach?]
Dla ustalonego $(t,x)$ wartość rozwiązania
zależy od $g,h$ (i pochodnych $g$) na sferze
$\{y\in\R^3:|y-x|=ct\}$. Dane z wnętrza kuli,
ale z dala od tej sfery, nie mają wpływu na
$u(t,x)$. Jest to silna zasada Huygensa dla
równania falowego w wymiarze trzy.
\end{uwaga}

\section{Dwa wymiary i zasada Huygensa}
\label{sec:pde5-huygens}

\begin{twierdzenie}[Wzór Poissona w wymiarze dwa]
\label{thm:pde5-2d}
Niech $c>0$, $g\in C_c^3(\R^2)$, $h\in C_c^2(\R^2)$.
Jedyne rozwiązanie $C^2$ równania falowego w
$\R^2$ z tymi danymi jest dane wzorem
\begin{equation}\label{eq:pde5-2d}
\begin{aligned}
 u(t,x)={}&\frac{\mathrm d}{\mathrm dt}
 \left[\frac1{2\pi c}\int_{|y-x|<ct}
 \frac{g(y)}{\sqrt{c^2t^2-|y-x|^2}}\,
 \mathrm d\lambda_2(y)\right]\\
 &+\frac1{2\pi c}\int_{|y-x|<ct}
 \frac{h(y)}{\sqrt{c^2t^2-|y-x|^2}}\,
 \mathrm d\lambda_2(y),\qquad t>0.
\end{aligned}
\end{equation}
\end{twierdzenie}
\begin{proof}
Rozszerzamy $g,h$ na $\R^3$ tak, aby były niezależne
od trzeciej zmiennej: $\widetilde g(y,z)=g(y)$,
$\widetilde h(y,z)=h(y)$. Wzór Kirchhoffa definiuje
wtedy rozwiązanie trójwymiarowe niezależne od $z$
(jednoznaczność ze wniosku~\ref{cor:pde5-uniqueness}).
Dla $r>0$ parametryzujemy dwie półsfery
$|y-x|^2+z^2=r^2$ przez
$z=\pm\sqrt{r^2-|y-x|^2}$. Element pola na każdej
półsferze wynosi
$r(r^2-|y-x|^2)^{-1/2}\,\mathrm d\lambda_2(y)$.
Dlatego dla $\varphi$ niezależnej od $z$ mamy
\[
 M_r\varphi(x,0)=\frac1{2\pi r}
 \int_{|y-x|<r}\frac{\varphi(y)}{\sqrt{r^2-|y-x|^2}}
 \,\mathrm d\lambda_2(y).
\]
Wstawiając $r=ct$ do wzoru Kirchhoffa,
otrzymujemy dokładnie \eqref{eq:pde5-2d}.
\end{proof}

Zwróćmy uwagę na różnicę geometryczną: w wymiarze
trzy całkowanie odbywa się po \emph{sferze},
a w wymiarze dwa po \emph{całym dysku} z osobliwą,
ale całkowalną wagą. Dla $g=0$ i nieujemnej,
niezerowej funkcji $h\in C_c^2(\R^2)$ druga całka
jest ściśle dodatnia, gdy dysk $B(x,ct)$ obejmuje
nośnik $h$ (przy $t$ odpowiednio dużym).
Oznacza to ogon fali, którego nie ma w wersji
trójwymiarowej.

\begin{uwaga}[Silna i słaba zasada Huygensa]
Słaba zasada oznacza, że wpływ danych początkowych
nie może wyjść poza domknięty stożek $|x-y|\le ct$.
W przedstawionym modelu jest to wniosek z
\ref{thm:pde5-cone}. Silna zasada oznacza zależność
od danych wyłącznie na czole stożka, czyli na
$|x-y|=ct$. Dla stałych współczynników w przestrzeni
Euklidesowej silna zasada zachodzi w wymiarach
nieparzystych $n\ge3$, lecz zawodzi w wymiarach
parzystych i w wymiarze jeden. Uogólnienia
na wszystkie wymiary wymagają dalszych rachunków
z rozwiązaniami fundamentalnymi; na tym etapie
udowodniliśmy przypadki $n=1,2,3$.
\end{uwaga}

\section{Zewnętrzne wymuszenie i wzór Duhamela dla fal}
\label{sec:pde5-forcing}

\begin{twierdzenie}[Wzór dla niejednorodnego równania w jednym wymiarze]
\label{thm:pde5-forcing}
Niech $c,T>0$, $g\in C^2(\R)$, $h\in C^1(\R)$,
$f\in C^1([0,T]\times\R)$ i załóżmy, że wszystkie
te funkcje oraz pochodne występujące we wzorze są
ciągłe (nie wymagamy ograniczoności w nieskończoności).
Wówczas funkcja
\begin{equation}\label{eq:pde5-forcing}
\begin{aligned}
 u(t,x)={}&\frac{g(x-ct)+g(x+ct)}2
 +\frac1{2c}\int_{x-ct}^{x+ct}h(y)\,\mathrm d\lambda_1(y)\\
 &+\frac1{2c}\int_0^t\int_{x-c(t-s)}^{x+c(t-s)}
 f(s,y)\,\mathrm d\lambda_1(y)\,\mathrm ds
\end{aligned}
\end{equation}
rozwiązuje $u_{tt}-c^2u_{xx}=f$, $u(0)=g$,
$u_t(0)=h$ na $[0,T]\times\R$. Jest jedyna
w klasie rozwiązań $C^2$.
\end{twierdzenie}
\begin{proof}
Jednorodna część rozwiązuje problem bez wymuszenia.
Oznaczmy trzeci składnik przez $v(t,x)$.
Różniczkując dwukrotnie z wykorzystaniem
reguły Leibniza otrzymujemy
\[
 v_t(t,x)=\frac12\int_0^t\bigl[
 f(s,x+c(t-s))+f(s,x-c(t-s))\bigr] \,\mathrm ds
\]
oraz
\[
 v_{tt}(t,x)=f(t,x)+\frac c2\int_0^t
 \bigl[f_x(s,x+c(t-s))-f_x(s,x-c(t-s))\bigr]
 \,\mathrm ds.
\]
Jednocześnie
$v_{xx}(t,x)=\frac1{2c}\int_0^t
 \bigl[f_x(s,x+c(t-s))-f_x(s,x-c(t-s))\bigr]
 \,\mathrm ds$.
Stąd $v_{tt}-c^2v_{xx}=f$ i
$v(0,x)=v_t(0,x)=0$. Jednoznaczność wynika
z wniosku~\ref{cor:pde5-uniqueness} zastosowanego
do różnicy rozwiązań.
\end{proof}

\section{Warunki brzegowe: struna na odcinku}
\label{sec:pde5-string}

Na odcinku $\Omega=(0,L)$ rozważmy
$u_{tt}-c^2u_{xx}=0$ z warunkiem
$u(t,0)=u(t,L)=0$. Zagadnienie można rozwiązać
metodą rozdzielenia zmiennych. Dla $m\in\N$
połóżmy $\omega_m=cm\pi/L$; wtedy
\[
 u_m(t,x)=\bigl(a_m\cos(\omega_mt)
 +b_m\sin(\omega_mt)\bigr)\sin(m\pi x/L)
\]
jest rozwiązaniem klasycznym. Jeśli zadane dane
początkowe są skończonymi kombinacjami funkcji
$\sin(m\pi x/L)$, dobieramy
$a_m$ jako współczynniki danych $g$,
a $b_m$ jako współczynniki $h$ podzielone przez
$\omega_m$.

\begin{przyklad}[Brak tłumienia modów]
Dla $g(x)=\sin(m\pi x/L)$ i $h=0$ dostajemy
$u(t,x)=\cos(\omega_mt)\sin(m\pi x/L)$.
Amplituda nie maleje monotonicznie, podczas gdy
odpowiedni mod równania ciepła zanika jak
$e^{-\kappa(m\pi/L)^2t}$. Właśnie dlatego
regularizujący charakter półgrupy cieplnej
nie przenosi się na grupę ewolucji falowej.
\end{przyklad}

\begin{stwierdzenie}[Energia struny]
\label{prop:pde5-stringenergy}
Jeśli $u\in C^2([0,T]\times[0,L])$ spełnia
$u_{tt}=c^2u_{xx}$ dla $0<t<T$, $0<x<L$,
a także $u(t,0)=u(t,L)=0$, to
\[
 E(t)=\frac12\int_0^L(u_t^2+c^2u_x^2)\,
 \mathrm d\lambda_1(x)
\]
jest stała na $[0,T]$.
\end{stwierdzenie}
\begin{proof}
Z warunku Dirichleta i regularności mamy
$u_t(t,0)=u_t(t,L)=0$ dla $0<t<T$. Zatem
$E'(t)=c^2[u_tu_x]_{x=0}^{x=L}=0$.
Ciągłość energii daje równość także w chwilach
$t=0,T$.
\end{proof}

\section{Dlaczego schemat numeryczny musi respektować stożek?}
\label{sec:pde5-numerics}

Na regularnej siatce $x_j=j\Delta x$, $t_m=m\Delta t$
przybliżamy drugie pochodne różnicami centralnymi:
\begin{equation}\label{eq:pde5-leapfrog}
 \frac{U_j^{m+1}-2U_j^m+U_j^{m-1}}{(\Delta t)^2}
 =c^2\frac{U_{j+1}^m-2U_j^m+U_{j-1}^m}{(\Delta x)^2}.
\end{equation}
Wprowadzamy liczbę Couranta $\lambda=c\Delta t/\Delta x$.
Dla modu $U_j^m=e^{i(kj\Delta x-\omega m\Delta t)}$
otrzymujemy relację dyspersji
\begin{equation}\label{eq:pde5-cfl}
 \sin^2\!\left(\frac{\omega\Delta t}2\right)
 =\lambda^2\sin^2\!\left(\frac{k\Delta x}2\right).
\end{equation}
Jeżeli $\lambda\le1$, istnieją rzeczywiste
częstotliwości $\omega$ dla wszystkich częstości
siatkowych. Gdy $\lambda>1$, istnieją częstości,
dla których prawa strona przekracza $1$, co
prowadzi do wykładniczego wzrostu rozwiązań
rekurencji. Jest to konieczny warunek stabilności
metody centralnej; w tym konkretnym jednowymiarowym
schemacie z prawidłowo zadanymi danymi
$\lambda\le1$ jest również klasycznym warunkiem
stabilności w odpowiedniej normie energetycznej
(z dodatkową ostrożnością dla częstości granicznych
przy $\lambda=1$). Pełny dowód numeryczny
pozostawimy do działu metod różnicowych.

\begin{uwaga}[Granica stabilności przy $\lambda=1$]
Dla częstotliwości siatkowej $k\Delta x=\pi$ amplituda
modu spełnia rekurencję $A_{m+1}+2A_m+A_{m-1}=0$.
Jej ogólne rozwiązanie to $(\alpha+\beta m)(-1)^m$:
podwójny pierwiastek $-1$ dopuszcza więc liniowy wzrost
przy dowolnie zadanych dwóch pierwszych warstwach.
Dla prędkości początkowej równej zeru i poprawnego
pierwszego kroku $A_1=-A_0$, skąd $\beta=0$.
Analiza stabilności w przypadku granicznym musi uwzględniać
przestrzeń danych i sposób inicjalizacji schematu;
proste kryterium rzeczywistości $\omega$ nie wystarcza.
\end{uwaga}

\begin{uwaga}[Zasięg numeryczny a zasięg fizyczny]
Schemat \eqref{eq:pde5-leapfrog} korzysta z
sąsiadów $j-1,j,j+1$ w poprzedniej warstwie
czasowej. Jeden krok może więc przenieść informację
o odległość $\Delta x$, podczas gdy fizyczna fala
przebywa $c\Delta t$. Warunek $c\Delta t\le\Delta x$
pozwala fizycznemu stożkowi zmieścić się w
numerycznym obszarze zależności. To tłumaczy
geometrię warunku CFL, ale nie zastępuje
dowodu stabilności.
\end{uwaga}

\section{Podsumowanie}
\label{sec:pde5-summary}

Dla równania falowego rozwiązanie w punkcie jest
wyznaczone przez dane w skończonym obszarze
początkowym; wynika to z kurczącej się energii
stożkowej. Równania ciepła i falowe są pod tym
względem fundamentalnie odmienne. Wzory d'Alemberta,
Kirchhoffa i Poissona pokazują również, że w samym
równaniu falowym silna zasada Huygensa zależy
od parzystości wymiaru. Następny rozdział poświęcimy
systematycznemu wykorzystaniu transformaty i
szeregów Fouriera, łącząc otrzymane tu konstrukcje
z metodami spektralnymi.

% ROZDZIAŁ 6: do wstawienia przed rozdziałem z kodem numerycznym.
